% Options for packages loaded elsewhere
% Options for packages loaded elsewhere
\PassOptionsToPackage{unicode}{hyperref}
\PassOptionsToPackage{hyphens}{url}
\PassOptionsToPackage{dvipsnames,svgnames,x11names}{xcolor}
%
\documentclass[
  letterpaper,
  DIV=11,
  numbers=noendperiod]{scrartcl}
\usepackage{xcolor}
\usepackage{amsmath,amssymb}
\setcounter{secnumdepth}{5}
\usepackage{iftex}
\ifPDFTeX
  \usepackage[T1]{fontenc}
  \usepackage[utf8]{inputenc}
  \usepackage{textcomp} % provide euro and other symbols
\else % if luatex or xetex
  \usepackage{unicode-math} % this also loads fontspec
  \defaultfontfeatures{Scale=MatchLowercase}
  \defaultfontfeatures[\rmfamily]{Ligatures=TeX,Scale=1}
\fi
\usepackage{lmodern}
\ifPDFTeX\else
  % xetex/luatex font selection
\fi
% Use upquote if available, for straight quotes in verbatim environments
\IfFileExists{upquote.sty}{\usepackage{upquote}}{}
\IfFileExists{microtype.sty}{% use microtype if available
  \usepackage[]{microtype}
  \UseMicrotypeSet[protrusion]{basicmath} % disable protrusion for tt fonts
}{}
\makeatletter
\@ifundefined{KOMAClassName}{% if non-KOMA class
  \IfFileExists{parskip.sty}{%
    \usepackage{parskip}
  }{% else
    \setlength{\parindent}{0pt}
    \setlength{\parskip}{6pt plus 2pt minus 1pt}}
}{% if KOMA class
  \KOMAoptions{parskip=half}}
\makeatother
% Make \paragraph and \subparagraph free-standing
\makeatletter
\ifx\paragraph\undefined\else
  \let\oldparagraph\paragraph
  \renewcommand{\paragraph}{
    \@ifstar
      \xxxParagraphStar
      \xxxParagraphNoStar
  }
  \newcommand{\xxxParagraphStar}[1]{\oldparagraph*{#1}\mbox{}}
  \newcommand{\xxxParagraphNoStar}[1]{\oldparagraph{#1}\mbox{}}
\fi
\ifx\subparagraph\undefined\else
  \let\oldsubparagraph\subparagraph
  \renewcommand{\subparagraph}{
    \@ifstar
      \xxxSubParagraphStar
      \xxxSubParagraphNoStar
  }
  \newcommand{\xxxSubParagraphStar}[1]{\oldsubparagraph*{#1}\mbox{}}
  \newcommand{\xxxSubParagraphNoStar}[1]{\oldsubparagraph{#1}\mbox{}}
\fi
\makeatother


\usepackage{longtable,booktabs,array}
\usepackage{calc} % for calculating minipage widths
% Correct order of tables after \paragraph or \subparagraph
\usepackage{etoolbox}
\makeatletter
\patchcmd\longtable{\par}{\if@noskipsec\mbox{}\fi\par}{}{}
\makeatother
% Allow footnotes in longtable head/foot
\IfFileExists{footnotehyper.sty}{\usepackage{footnotehyper}}{\usepackage{footnote}}
\makesavenoteenv{longtable}
\usepackage{graphicx}
\makeatletter
\newsavebox\pandoc@box
\newcommand*\pandocbounded[1]{% scales image to fit in text height/width
  \sbox\pandoc@box{#1}%
  \Gscale@div\@tempa{\textheight}{\dimexpr\ht\pandoc@box+\dp\pandoc@box\relax}%
  \Gscale@div\@tempb{\linewidth}{\wd\pandoc@box}%
  \ifdim\@tempb\p@<\@tempa\p@\let\@tempa\@tempb\fi% select the smaller of both
  \ifdim\@tempa\p@<\p@\scalebox{\@tempa}{\usebox\pandoc@box}%
  \else\usebox{\pandoc@box}%
  \fi%
}
% Set default figure placement to htbp
\def\fps@figure{htbp}
\makeatother





\setlength{\emergencystretch}{3em} % prevent overfull lines

\providecommand{\tightlist}{%
  \setlength{\itemsep}{0pt}\setlength{\parskip}{0pt}}



 


\KOMAoption{captions}{tableheading}
\makeatletter
\@ifpackageloaded{caption}{}{\usepackage{caption}}
\AtBeginDocument{%
\ifdefined\contentsname
  \renewcommand*\contentsname{Table of contents}
\else
  \newcommand\contentsname{Table of contents}
\fi
\ifdefined\listfigurename
  \renewcommand*\listfigurename{List of Figures}
\else
  \newcommand\listfigurename{List of Figures}
\fi
\ifdefined\listtablename
  \renewcommand*\listtablename{List of Tables}
\else
  \newcommand\listtablename{List of Tables}
\fi
\ifdefined\figurename
  \renewcommand*\figurename{Figure}
\else
  \newcommand\figurename{Figure}
\fi
\ifdefined\tablename
  \renewcommand*\tablename{Table}
\else
  \newcommand\tablename{Table}
\fi
}
\@ifpackageloaded{float}{}{\usepackage{float}}
\floatstyle{ruled}
\@ifundefined{c@chapter}{\newfloat{codelisting}{h}{lop}}{\newfloat{codelisting}{h}{lop}[chapter]}
\floatname{codelisting}{Listing}
\newcommand*\listoflistings{\listof{codelisting}{List of Listings}}
\makeatother
\makeatletter
\makeatother
\makeatletter
\@ifpackageloaded{caption}{}{\usepackage{caption}}
\@ifpackageloaded{subcaption}{}{\usepackage{subcaption}}
\makeatother
\usepackage{bookmark}
\IfFileExists{xurl.sty}{\usepackage{xurl}}{} % add URL line breaks if available
\urlstyle{same}
\hypersetup{
  pdftitle={Lecture 2: Measuring Yield},
  colorlinks=true,
  linkcolor={blue},
  filecolor={Maroon},
  citecolor={Blue},
  urlcolor={Blue},
  pdfcreator={LaTeX via pandoc}}


\title{Lecture 2: Measuring Yield}
\author{}
\date{}
\begin{document}
\maketitle

\renewcommand*\contentsname{Table of contents}
{
\hypersetup{linkcolor=}
\setcounter{tocdepth}{3}
\tableofcontents
}

\section{Learning Objectives}\label{learning-objectives}

\begin{itemize}
\tightlist
\item
  Compute the yield (internal rate of return) on any investment using
  the present‑value equation that equates the price of the investment to
  the present value of its cash‑flows【965351366966888†L3643-L3669】.
\item
  Calculate conventional yield measures---current yield, yield to
  maturity, yield to call, yield to put, cash‑flow yield---and
  understand how each is interpreted【965351366966888†L3896-L3973】.
\item
  Calculate the yield for a portfolio by discounting all portfolio
  cash‑flows so that the present value equals the portfolio price.
\item
  Compute the discount margin for a floating‑rate security and recognise
  why this measure differs from standard yield measures.
\item
  Identify the three potential sources of a bond's return (coupon
  interest, reinvestment income and capital gain or loss) and explain
  the concept of reinvestment risk.
\item
  Recognise the limitations of conventional yield measures and why total
  return is a more realistic performance metric.
\item
  Calculate the total return on a bond for a specified investment
  horizon under explicit reinvestment rate assumptions and a projected
  sale price.
\item
  Use horizon analysis to assess how reinvestment assumptions and future
  yield expectations influence bond returns.
\item
  Measure yield changes using both absolute (basis‑point) and relative
  (percentage) conventions.
\end{itemize}

\begin{center}\rule{0.5\linewidth}{0.5pt}\end{center}

\section{1~Computing the Yield or Internal Rate of
Return}\label{computing-the-yield-or-internal-rate-of-return}

\subsection{1.1~Yield on Any Investment}\label{yield-on-any-investment}

The \textbf{yield} or \textbf{internal rate of return (IRR)} on an
investment is the interest rate that makes the present value of the
investment's cash‑flows equal to its price. Mathematically, for an
investment with (N) future cash‑flows (CF\_t) and price (P), the yield
(y) satisfies

{[} P ;=; \sum\_\{t=1\}\^{}\{N\}\frac{CF_t}{(1+y)^t}. {]}

Solving for (y) generally requires an iterative procedure because (y)
appears in the denominator of each term【965351366966888†L3643-L3669】.
A higher assumed yield lowers the present value; a lower assumed yield
raises it. The objective is to find the yield for which the present
value equals the price.

\subsubsection{Special case: a single cash
flow}\label{special-case-a-single-cash-flow}

If the investment has just one future cash flow (CF\_N) in (N) periods,
the yield can be solved algebraically: {[} y ;=;
\left(\frac{CF_N}{P}\right)\^{}\{1/N\} - 1. {]}

This formula eliminates the need for trial and error because there is
only one cash‑flow term【965351366966888†L3643-L3669】.

\subsection{1.2~Example: Computing Yield by Trial and
Error}\label{example-computing-yield-by-trial-and-error}

Suppose a financial instrument with price \$903.10 promises cash‑flows
of \$100, \$100, \$100 and \$1,000 in years~1--4. To find its yield:

\begin{enumerate}
\def\labelenumi{\arabic{enumi}.}
\tightlist
\item
  Compute the present value of cash‑flows at several interest rates.
\item
  Adjust the rate until the present value matches \$903.10.
\end{enumerate}

Using a 10~\% trial rate gives a present value of \$931.69. A 12~\%
trial rate produces \$875.71. By interpolation, a rate of 11~\% equates
the present value to the price【965351366966888†L3718-L3767】. The yield
on this investment is therefore \textbf{11~\%}.

\subsection{1.3~Periodic Cash Flows and
Annualisation}\label{periodic-cash-flows-and-annualisation}

Cash‑flows are often paid more than once per year---for example,
semiannual coupons. If there are (n) compounding periods per year, the
yield is the rate per period ((y)) that satisfies {[} P =
\sum\_\{t=1\}\^{}\{nN\} \frac{CF_t}{(1+y)^t}. {]}

For semiannual coupons, doubling (y) gives the \textbf{bond‑equivalent
yield}; for other frequencies, multiply (y) by the number of periods in
the year. To find the \textbf{effective annual yield} corresponding to a
periodic rate (r\_p), compute {[} \text{Effective annual yield} ;=; (1 +
r\_p)\^{}m - 1, {]} where (m) is the number of periods per year. For
example, a quarterly periodic rate of 2~\% implies an effective annual
yield of ((1.02)\^{}4-1\approx8.24\%).

\begin{center}\rule{0.5\linewidth}{0.5pt}\end{center}

\section{2~Conventional Yield
Measures}\label{conventional-yield-measures}

Bonds are commonly quoted using several yield measures that summarise
different aspects of return. Each measure uses the basic present‑value
relationship but makes different assumptions about cash‑flows and
reinvestment.

\subsection{2.1~Current Yield}\label{current-yield}

The \textbf{current yield} relates a bond's annual coupon to its market
price: {[} \text{current yield} ;=;
\frac{\text{annual coupon}}{\text{price}}. {]}

For instance, a 15‑year bond with a 7~\% coupon and price \$769.42 has a
current yield of (70/769.42 \approx 9.10\%). A 4.5~\% coupon bond priced
at 99.531 (per \$100 par) has a current yield of (4.50/99.531
\approx 4.52\%)【965351366966888†L3896-L3922】.

The current yield ignores capital gains or losses and the time value of
money. It only measures coupon income relative to price.

\subsection{2.2~Yield to Maturity (YTM)}\label{yield-to-maturity-ytm}

The \textbf{yield to maturity} is the internal rate of return assuming
the bond is held to maturity and all coupon payments are reinvested at
the same rate. For a bond with price (P), coupon payment (C), maturity
value (M) and (n) semiannual periods, the present‑value equation is {[}
P ;=; \sum\_\{t=1\}\^{}\{n\}\frac{C}{(1+y)^t} ;+; \frac{M}{(1+y)^n}. {]}

Solving for (y) (a semiannual rate) generally requires trial and error;
doubling (y) gives the bond‑equivalent
YTM【965351366966888†L3896-L3973】. YTM incorporates coupon income,
reinvestment income and any capital gain or loss realised at maturity.

\subsection{2.3~Yield to Call, Yield to Put and Yield to
Sinker}\label{yield-to-call-yield-to-put-and-yield-to-sinker}

\textbf{Callable bonds} allow the issuer to redeem the bond at specified
call dates for a predetermined price. The \textbf{yield to call} is the
interest rate that equates the bond's price to the present value of the
coupon payments up to the assumed call date plus the call price. Because
multiple call dates may exist, investors often compute yields to first
call, next call and to any refunding call.

Similarly, \textbf{putable bonds} allow the investor to sell the bond
back to the issuer on specified dates at a put price. The \textbf{yield
to put} equates the price to the present value of cash‑flows up to the
put date plus the put price. Bonds with sinking‑fund schedules retire
portions of the issue over time. The \textbf{yield to sinker} uses the
sinking‑fund date and price.

The \textbf{yield to worst} is the lowest yield among the yield to
maturity, all yields to call and all yields to put. It represents the
minimum return an investor might realise if the issuer exercises
embedded options.

\subsection{2.4~Yield for a Portfolio}\label{yield-for-a-portfolio}

The yield of a bond portfolio is not simply the average of the
individual yields because each bond has a different pattern of
cash‑flows. To compute a portfolio yield, aggregate the portfolio's
cash‑flows and find the discount rate that makes the present value of
all cash‑flows equal the portfolio's price. For instance, a portfolio of
three bonds with differing maturities and coupon rates requires
discounting the combined cash‑flow schedule to match the total market
value.

\subsection{2.5~Cash‑Flow Yield and Discount
Margin}\label{cashflow-yield-and-discount-margin}

For amortising securities such as mortgage‑backed or asset‑backed
securities, investors receive principal repayments along with interest.
The \textbf{cash‑flow yield} is the rate that equates the price to the
projected cash‑flows, including scheduled and unscheduled principal
repayments. Estimating the cash‑flow yield requires assumptions about
prepayments and reinvestment.

Floating‑rate securities pay interest linked to a reference rate (e.g.,
LIBOR) plus a quoted margin. The cash‑flows depend on future reference
rates, so a YTM cannot be calculated directly. Investors instead
calculate the \textbf{discount margin}, the margin over the reference
rate that equates the price to the present value of the projected
cash‑flows assuming the reference rate remains constant. This measure
estimates the average spread an investor can expect over the reference
rate.

\begin{center}\rule{0.5\linewidth}{0.5pt}\end{center}

\section{3~Sources of a Bond's Return and Reinvestment
Risk}\label{sources-of-a-bonds-return-and-reinvestment-risk}

A bond's total dollar return comes from three sources:

\begin{enumerate}
\def\labelenumi{\arabic{enumi}.}
\tightlist
\item
  \textbf{Coupon interest} -- the periodic coupon payments.
\item
  \textbf{Capital gain or loss} -- the difference between the sale price
  (or maturity value) and the purchase price.
\item
  \textbf{Reinvestment income} -- interest earned from reinvesting the
  coupon payments (and, for amortising bonds, reinvested principal).
\end{enumerate}

Conventional yields implicitly assume that coupon payments can be
reinvested at the computed yield. The risk that future reinvestment
rates will be lower than the yield at purchase is \textbf{reinvestment
risk}. The longer the maturity and the higher the coupon, the more the
bond's total return depends on the reinvestment of coupon payments, and
therefore the greater the reinvestment risk.

\subsection{3.1~Interest‐on‐Interest
Component}\label{interestoninterest-component}

For a non‑amortising bond paying semiannual coupon (C) over (n) periods
and reinvestment rate (r), the future value of the coupon payments plus
reinvestment income after (n) periods is {[} C,\frac{(1+r)^n - 1}{r}.
{]}

The interest‑on‑interest component is the difference between this amount
and the sum of the coupon payments ((nC)). For long‑maturity,
high‑coupon bonds, the interest‑on‑interest component may represent a
significant portion of total return. Zero‑coupon bonds have no
reinvestment risk because all return is realised at maturity.

\begin{center}\rule{0.5\linewidth}{0.5pt}\end{center}

\section{4~Total Return and Horizon
Analysis}\label{total-return-and-horizon-analysis}

The yield to maturity and other conventional yields are \textbf{promised
yields}: they are realised only if the bond is held to the date used in
the calculation and coupon payments can be reinvested at that yield. To
evaluate a bond's performance over a specific \textbf{investment
horizon}, investors compute the \textbf{total return}, which explicitly
incorporates assumptions about reinvestment rates and the bond's sale
price at the end of the horizon.

\subsection{4.1~Computing Total Return}\label{computing-total-return}

To compute the total return over a horizon of (h) semiannual periods:

\begin{enumerate}
\def\labelenumi{\arabic{enumi}.}
\tightlist
\item
  \textbf{Reinvestment income} -- For each coupon payment received
  before the horizon, compute its future value at the assumed
  reinvestment rate until the horizon. Summing these gives the future
  value of coupon payments plus reinvestment income.
\item
  \textbf{Projected sale price} -- Estimate the bond's price at the
  horizon by discounting its remaining cash‑flows at the expected yield
  to maturity prevailing at that time.
\item
  \textbf{Total future dollars} -- Add the future value of coupon
  payments to the projected sale price. Subtract the initial purchase
  price to obtain the dollar return.
\item
  \textbf{Total return} -- Solve for the interest rate that equates the
  purchase price to the total future dollars; double the semiannual rate
  to express as an annualised return.
\end{enumerate}

This procedure accommodates multiple reinvestment rates. For example,
coupon payments received in early periods may be reinvested at one rate
and later payments at another. The final projected price depends on
expectations for the yield environment at the horizon.

\subsection{4.2~Example: Total Return}\label{example-total-return}

An investor with a three‑year horizon considers a 20‑year 8~\% bond
priced at \$828.40 with YTM 10~\%. The investor expects to reinvest
coupons at 6~\% and to sell the bond after three years at a yield of
7~\% for then‑17‑year bonds. The calculation proceeds as follows:

\begin{enumerate}
\def\labelenumi{\arabic{enumi}.}
\tightlist
\item
  The bond pays \$40 every six months. Assuming 3~\% semiannual
  reinvestment, the future value of coupon payments over six periods is
  \$258.74.
\item
  The projected sale price in three years (with a 7~\% yield) is
  \$1\,098.51.
\item
  Total future dollars are \$1\,357.25; dividing by the purchase price
  and solving for the semiannual return gives 8.58~\% per half year.
\item
  Doubling yields an \textbf{annualised total return of 17.16~\%},
  demonstrating how reinvestment and expected sale price affect realised
  performance.
\end{enumerate}

Horizon analysis is also used in bond‑swap decisions. By calculating
total returns under different reinvestment and yield scenarios,
investors can determine whether swapping one bond for another improves
projected performance.

\begin{center}\rule{0.5\linewidth}{0.5pt}\end{center}

\section{5~Measuring Yield Changes}\label{measuring-yield-changes}

Yield changes can be measured in two ways:

\begin{itemize}
\item
  \textbf{Absolute yield change} -- measured in basis points. It is the
  absolute difference between two yields: {[}
  \text{absolute change (bp)} ;=; \textbar y\_\{\text{new}\} -
  y\_\{\text{old}\}\textbar{} \times 100. {]}
\item
  \textbf{Relative (percentage) yield change} -- calculated as the
  natural logarithm of the ratio of the new yield to the old yield: {[}
  \text{relative change (\%)} ;=; 100
  \times \ln\left(\frac{y_{\text{new}}}{y_{\text{old}}}\right). {]}
\end{itemize}

For example, if yields rise from 4.45~\% to 5.11~\%, the absolute change
is 66~bp and the relative change is (\ln(5.11/4.45)\times 100
\approx 13.83\%). These measures are useful for comparing yield
movements across different sectors and maturities.

\begin{center}\rule{0.5\linewidth}{0.5pt}\end{center}

\section{6~Key Points}\label{key-points}

\begin{itemize}
\tightlist
\item
  The yield on any investment is the internal rate of return that
  equates the present value of cash‑flows to the
  price【965351366966888†L3643-L3669】.
\item
  Conventional yield measures (current yield, yield to maturity, yield
  to call, etc.) summarise different aspects of return and rely on
  specific reinvestment and holding‑period
  assumptions【965351366966888†L3896-L3973】.
\item
  A bond's total dollar return consists of coupon interest, reinvestment
  income and any capital gain or loss; reinvestment risk arises because
  future rates are uncertain.
\item
  Total return is a more comprehensive measure than yield to maturity
  because it incorporates explicit reinvestment assumptions and a
  horizon‑based sale price.
\item
  Horizon analysis uses the total return framework to compare bonds or
  strategies over a specified holding period under alternative
  interest‑rate scenarios.
\end{itemize}




\end{document}
